\subsection{Scope}

The complete recipe is better than the prior arm on all four fleets, but the size of the gain tracks the number of operating conditions. On the two single-condition fleets it is 0.79 and 2.40 percent better (FD001 10.154 against 10.235, FD003 9.914 against 10.158); on the two six-condition fleets it is 10.00 and 9.44 percent better.

Mixup alone shows the opposite pattern. On the single-condition fleets it is 1.04 and 4.67 percent worse than the prior arm, while on FD003 it remains 2.13 percent better than the plain backbone. The prior therefore carries information that nothing else supplies when the operating condition is fixed; the span schedule recovers that loss, and under six conditions the mixed-condition regularisation is worth more than the prior in the first place.

We read the pattern as a scope statement rather than a caveat: the schedule exploits the longer history where that history is more informative, which is under varying conditions, and it still pays a small dividend when the condition is fixed.
